\documentclass[9pt,aspectratio=169]{beamer}

\usepackage[T1,T2A]{fontenc}

\usepackage{pgfplots}
\usepackage{pgfplotstable}
\pgfplotsset{compat=1.18}
\usepgfplotslibrary{statistics}

\usepackage{caption}

\usepackage[
    backend=biber,
    style=gost-numeric,
    language=auto,
    bibencoding=utf8,
    url=false,
    doi=false,
    isbn=false,
    defernumbers=true,
]{biblatex}

%\addbibresource{references.bib}

\defbibfilter{englishonly}{
  keyword=english
}

\defbibfilter{russianonly}{
  keyword=russian
}

\defbibfilter{links}{
  keyword=link
}

\usepackage{makecell}

\usepackage[utf8]{inputenc}
\usepackage{amsfonts}
\usepackage{amssymb}
\usepackage{float}
\usepackage{amsmath}
\usepackage[english,russian]{babel}

%  - Гиперссылки (url, \ref-ссылки, \cite-ссылки)
\usepackage{hyperref}
\usepackage{ragged2e}

\usetheme{Madrid}

\setbeamersize{text margin left=0.9cm, text margin right=0.9cm}

\setbeamertemplate{navigation symbols}{}
\setbeamertemplate{footline}[frame number]

\setbeamertemplate{title page}{
    \begin{centering}
    
    \vspace{1cm}

    \begin{beamercolorbox}[sep=8pt,center,rounded=true,shadow=true]{title}
        \usebeamerfont{title}\inserttitle
    \end{beamercolorbox}
    \vspace{0.8cm}

    {\usebeamerfont{author}\insertauthor\par}
    \vspace{0.8cm}

    \begin{minipage}{0.6\textwidth}
    \end{minipage}
    \vfill
    \hfill
    \begin{minipage}{0.31\textwidth}
        \raggedright
        Научный руководитель: \\
        \;\;\;\;\;к.ф.-м.н., доцент \\
        \;\;\;\;\;Белянкин Г.А.
    \end{minipage}

    \vfill
    \vfill
    \vfill
    \vfill
    \vfill
    {\usebeamerfont{date}\insertdate\par}

    \end{centering}
}
\setbeamerfont{title}{size=\LARGE,series=\bfseries}

\setbeamerfont{frametitle}{size=\LARGE,series=\bfseries}

\setbeamertemplate{enumerate items}[square]

\title{Реконструкция относительной трехмерной позы человека
по необработанным видеоданным в условиях отсутствия метаданных}
\author{Фостенко Олег Андреевич, группа 411 \\
Факультет вычислительной математики и кибернетики \\
МГУ~им.~М.~В.~Ломоносова, кафедра исследования операций}
\date{Москва, 2026}

\begin{document}

\frame{\titlepage}





\begin{frame}{Актуальность и цель}

Оценка (восстановление) положения ключевых точек тела человека в пространстве важна
для количественного анализа движений.

Большая часть современных подходов требует предварительной калибровки
камер и использования специализированного
оборудования.

В данной работе предлагается метод восстановления трехмерной позы, не
требующий настройки и предъявляющий минимальные требования к оборудованию.

\vfill 

\begin{block}{Цель}
Разработать метод восстановления трехмерной позы человека
по необработанным видеоданным без калибровки камер при наличии в качестве
оборудования лишь одной или нескольких (двух) камер.
\end{block}


\end{frame}






\begin{frame}{Цель}

%REMOVE AFTER REMOVE AFTER REMOVE AFTER

\begin{columns}

\column{0.5\textwidth}
\centering
<Блок с изображением некоторого кадра из видео>

\column{0.5\textwidth}
\centering
<Блок с изображением трехмерных ключевых точек
(визуализация Matplotlib, соответствующая кадру слева на данном слайде)>

\end{columns}
%REMOVE BEFORE REMOVE BEFORE REMOVE BEFORE

\end{frame}





\begin{frame}{Задачи}

\begin{columns}

% LEFT SIDE: TEXT
\column{0.65\textwidth}

Вход:
видеоряд \(V = (I_1, \dots, I_T), \quad \text{где } I_t \in \mathbb{R}^{H \times W \times 3}\) --- $t$-й кадр.

\vspace{0.5em}

Требуется построить способ обработки видео, включающий следующие этапы:

\begin{enumerate}
    \justifying
    \item Оценка 2D позы (ключевые точки на видео):
    \[
        S = (S_1,\dots,S_T), \quad S_t \text{ --- 2D поза на $t$-м кадре:}
    \]
    \[
        S_t = (p_t^1(I_t), \dots, p_t^K(I_t)), \quad p_t^j = (u_t^j,v_t^j,c_t^j) \in \mathbb{R}^3,
    \]
    где $p_t^j$ --- тройка, характеризующая двумерную ключевую точку, в которой
    \((u^t_j,v^t_j)\) --- координаты точки в пиксельном пространстве на $t$-м кадре,
    \(c_j\) --- уверенность модели на $t$-м кадре,
    \(K\) --- число ключевых точек, описывающих 2D позу человека, \(t=1,\dots,T\).

\end{enumerate}

% RIGHT SIDE: IMAGE
\column{0.35\textwidth}
<Блок с изображением двумерной скелетной модели, наложенной
на некоторый кадр из видео>

\end{columns}

\end{frame}



\begin{frame}{Задачи}

\begin{columns}

% LEFT SIDE: TEXT
\column{0.65\textwidth}

\begin{enumerate}
\setcounter{enumi}{1}
    \item Восстановление относительной трехмерной позы:
    \[
        P = (P_1, \dots, P_T), \quad P_t \in \mathbb{R}^{3 \times M},
    \]
    где \(M\) -- число ключевых точек, описывающих 3D позу человека, \(t=1,\dots,T\).
\end{enumerate}


% RIGHT SIDE: IMAGE
\column{0.35\textwidth}
<Блок с изображением трехмерной скелетной модели>

\end{columns}

\end{frame}







\begin{frame}{Задачи}

\begin{columns}

% LEFT SIDE: TEXT
\column{0.65\textwidth}

Относительная трехмерная поза --- набор ключевых точек скелета
без глобального положения в сцене.

Глобальная трехмерная поза --- положение и ориентация
в системе координат камеры.

\begin{enumerate}
\setcounter{enumi}{2}
    \item Разработка метода оценки трехмерной позы в мировых
    координатах при наличии одной камеры (монокулярный режим)
    \item Разработка метода оценки трехмерной позы в мировых
    координатах при наличии двух камер (режим стерео).

\end{enumerate}


% RIGHT SIDE: IMAGE
\column{0.35\textwidth}
<Блок с изображением трехмерной скелетной модели>

\end{columns}

\end{frame}







\begin{frame}{Оценка относительной позы}

\begin{columns}

% LEFT SIDE: BLOCK_0
\column{0.65\textwidth}

Для оценки 2D позы и получения относительной 3D позы
используются нейросетевые модели.

\vspace{0.8em}

Для подавления артефактов на стыках окон при предсказании 
трехмерной позы используется
временное сглаживание с гауссовским ядром.

\vspace{0.2em}

<Диаграмма, графически показывающая перекрывающиеся окна с
визуализацией гауссовских "колоколов" (как здесь
\url{https://miro.medium.com/v2/resize:fit:1100/format:webp/
1*O0wfmYfbkrMZxkMaVqR7ww.png})\:
>

% RIGHT SIDE: IMAGE
\column{0.35\textwidth}
<
Простая диаграмма со следующим текстом:
"Оценка 2D позы покадрово с помощью предобученной сверточной
нейросети на базе CSPNeXt" 
$\rightarrow$ "Оценка 3D позы на окнах длиной 243 кадра
предобученной трансформерной модели на базе DSTFormer" 
$\rightarrow$ "Временное сглаживание"\:>


\end{columns}

\end{frame}








\begin{frame}{Основные подзадачи}


\end{frame}












\begin{block}{Задачи}
\begin{itemize}
    \item разработка метода восстановления относительной 3D позы человека по видеоданным;
    \item создание метода оценки глобальной позы в моно-камерном режиме (одна камера);
    \item создание метода оценки глобальной позы в режиме стерео (две камеры).
\end{itemize}
\end{block}

В данной работе основное внимание уделяется применению методов компьютерного
зрения для автоматического восстановления 2D- и 3D-координат ключевых
суставных точек человека на основе видео. Для решения задачи используются
state-of-the-art (SOTA) нейросетевые модели, обученные на больших
датасетах с разметкой ключевых точек.

В отличие от подходов, основанных на использовании носимых
устройств (сенсоров)~\cite{conf1,conf2}, предлагаемый метод не
требует дополнительного оборудования. В работе рассматриваются два
варианта съемки: моно- и мультикамерный (с двумя камерами). При этом
оба варианта реализованы без проведения калибровки камер и точного
определения их взаимного расположения. Это делает подход потенциально
применимым в ситуациях с ограниченным техническим оснащением.




\end{document}